\chapter{Linear Advance Calibration} \label{app:LinearAdvance}

To investigate whether extrusion consistency could be improved, the Marlin Linear Advance feature was evaluated. Linear Advance compensates for pressure variations in the nozzle during acceleration and deceleration by adjusting the extrusion rate based on the estimated nozzle pressure. According to the Marlin documentation, this can reduce over-extrusion and under-extrusion during changes in print speed and improve dimensional accuracy and line uniformity \cite{Lin_adv}.

A Linear Advance calibration print is shown in \cref{fig:Lin_adv}. During the test, the print speed alternates between 20~mm/s and 40~mm/s before returning to 20~mm/s. Each horizontal row corresponds to a different Linear Advance factor, $K$, ranging from 0 to 3 in increments of 0.1. The purpose of the test is to identify the $K$-value that produces the most consistent extrusion width despite the imposed speed changes.

The colored boxes in \cref{fig:Lin_adv} highlight representative examples of the observed extrusion behavior. The red boxes indicate under-extrusion, where the deposited trace becomes narrower than intended after a speed transition. The blue boxes indicate over-extrusion, where excess material results in a wider trace. The green boxes highlight regions where the extrusion width remains relatively constant through the speed transitions, indicating successful compensation of nozzle pressure effects.

Although several $K$-values produced visually improved results in the calibration print, no clear improvement was observed when the feature was applied to the custom-generated serpentine toolpaths used in this work. In some cases, the extrusion consistency was even degraded. Consequently, Linear Advance was not used in the final printing procedure.


\begin{figure}
    \centering
    \includegraphics[width=0.75\linewidth]{Pictures/Lin_adv.png}
    \caption{Linear Advance calibration print. The print speed alternates between 20~mm/s and 40~mm/s, while the Linear Advance factor $K$ increases from 0 to 3 in steps of 0.1 between horizontal rows. Red boxes indicate under-extrusion, blue boxes indicate over-extrusion, and green boxes indicate regions where the extrusion width remains relatively constant despite the speed changes.}
    \label{fig:Lin_adv}
\end{figure}